\chapter{Conclusions and Outlook}
\label{chap:conclusion}

This thesis takes the European option pricing problem under the 3/2 stochastic volatility model as its core research subject. Addressing the shortcomings of existing research—where performance evaluations of pricing methods are largely confined to idealized academic benchmark cases and lack empirical validation in real extreme market environments—this study systematically constructed a comprehensive theoretical framework for eleven European option pricing methods under the 3/2 model. Through multi-scenario numerical experiments, it quantitatively evaluated the pricing accuracy, computational efficiency, and numerical stability of different methods. Furthermore, based on the extreme BTC market flash crash event of October 2025, it completed a full-cycle parameter calibration of the model and an extreme-scenario re-evaluation of the pricing methods. This clarified the performance boundaries of different methods in both idealized environments and real markets, providing a systematic analytical framework and empirical evidence for the theoretical application and engineering practice of the 3/2 model in derivatives pricing. This chapter will systematically summarize the core research conclusions of the entire thesis, objectively analyze the limitations of the study, and provide an outlook on future related research directions.

\section{Main Research Conclusions}

Through systematic research across three dimensions—theoretical construction, numerical experiments, and empirical analysis—this thesis has formed four core research conclusions that comprehensively cover the pricing method performance, market adaptability, and practical application boundaries of the 3/2 model.

First, this thesis systematically clarified the theoretical logic and performance boundaries of eleven option pricing methods for the 3/2 model within a pure diffusion framework, forming a clear hierarchical system of method performance. This thesis categorized the eleven pricing methods into three major paradigms: time-stepping conditional Monte Carlo methods, terminal exact and near-exact simulation methods, and Fourier frequency-domain pricing methods, clarifying the implementation pathways and theoretical assumptions of each method. Based on 7 sets of academic benchmark cases covering conventional markets and extreme parameter environments, a systematic performance evaluation of these methods was completed. The research results indicate that the eleven methods can be classified into five tiers based on engineering practicality: The first tier consists of the Fourier-domain COS and FFT methods, which are the only methods that maintain numerical stability across all scenarios while possessing extremely high pricing accuracy and computational efficiency, with the COS method demonstrating comprehensive superiority. The second tier includes the time-stepping Exact Stepping, Poisson-Gamma, and QE Near-Exact Stepping methods; these are the only solutions within the Monte Carlo framework experiencing zero numerical failures across all scenarios, balancing pricing accuracy and computational efficiency, and serving as the sole reliable choices for scenarios requiring complete path generation. The third tier is the Euler/Milstein discretization scheme, which achieves acceptable pricing accuracy only at extremely fine time steps and has extremely low computational efficiency, serving merely as a benchmark tool for theoretical accuracy comparisons. The fourth tier comprises the IG and Gamma near-exact simulation methods, whose biases are controllable only in specific ATM scenarios but exhibit systematic and massive pricing deviations in the vast majority of scenarios, making them suitable only for academic comparisons in specific settings. The complete failure tier consists of the Baldeaux exact simulation, Fourier-Cosine exact terminal simulation, and moment matching exact terminal simulation methods, which can only output valid results in a single benchmark scenario and completely fail numerically in all other scenarios. The core finding is that the theoretical unbiasedness and ``exact'' label of a pricing method do not equate to full-scenario stability in engineering practice; most nominally ``exact'' terminal simulation methods can only complete theoretical validation under idealized benchmark parameters and fail to adapt to non-benchmark extreme parameter environments.

Second, the 3/2 stochastic volatility model can effectively capture option pricing characteristics in highly volatile cryptocurrency markets, demonstrating good market adaptability under extreme conditions. This thesis constructed a two-step calibration framework of ``long-term parameter time series estimation - time-varying parameter cross-sectional calibration.'' It estimated the long-term mean reversion parameters based on daily DVOL data from the half-year preceding the crash event and completed the cross-sectional calibration of time-varying parameters throughout the crash event cycle using the Medvedev-Scaillet asymptotic expansion method. Empirical results show that the 3/2 model under a pure diffusion framework can accurately fit the morphological evolution of the volatility smile across the crash event cycle. The model's core parameters exhibit significant event-driven time-varying patterns: ATM implied volatility and the volatility of volatility (VoV) spiked synchronously during the crash, with the increase for the 6-day short-term maturity being significantly higher than that for the 20-day medium-term maturity. This reflects a significant elevation in the market's short-term volatility expectations and tail risk pricing regarding extreme events. The negative correlation coefficient between price and volatility weakened in stages during the crash, with the distortion of the leverage effect being more pronounced for short-term option contracts, while the correlation for medium-term contracts remained relatively stable. This embodies the significant impact of panic trading and liquidity shocks on the short-term option pricing structure under extreme conditions. This result validates the good adaptability of the 3/2 model in highly volatile cryptocurrency markets, providing a reliable model foundation for pricing and risk measurement in crypto derivatives markets.

Third, real extreme market scenarios impose requirements on the numerical stability and robustness of pricing methods that are far higher than those of academic benchmark cases, further amplifying the performance divergence among different methods. Based on 6 sets of real market parameters obtained from calibration, this thesis constructed brand-new extreme scenario test cases, completed a performance re-evaluation of the pricing methods, and conducted a horizontal comparison with the academic benchmark results. The study found that terminal exact simulation methods, which only failed in a few scenarios in the academic benchmark cases, experienced full-scenario numerical failures or systematic and massive pricing biases in real extreme scenarios, completely losing their engineering practicality. Conversely, time-stepping distribution sampling methods and Fourier frequency-domain methods maintained 100\% numerical stability and stable pricing accuracy in both academic cases and real extreme scenarios, demonstrating extremely strong robustness against extreme conditions. Simultaneously, extreme scenarios further highlighted the performance advantages of the COS method, which maintained a pricing bias approaching 0 across all scenarios with a single-case computation time stabilizing around 0.05 seconds, proving to be the optimal scheme for batch pricing of European options, real-time valuation, and model calibration in extreme market environments. The Poisson-Gamma and QE near-exact stepping methods are the optimal choices for path-dependent option pricing and dynamic hedging scenarios, completely compensating for the core defects of the traditional Euler/Milstein schemes, which are sensitive to step size and exhibit extremely poor accuracy in extreme scenarios.

Fourth, based on theoretical and empirical results, this thesis has formed scenario-based application guidelines for option pricing methods under the 3/2 model. For scenarios such as the batch pricing of European options and model parameter calibration, the COS frequency-domain method should be prioritized; its full-scenario stability, supreme accuracy, and optimal efficiency fully satisfy the demands of these scenarios. For scenarios requiring complete variance path generation, such as path-dependent option pricing and dynamic hedging, the Poisson-Gamma or QE near-exact stepping methods should be prioritized; their balance of full-scenario stability, accuracy, and efficiency meets the engineering requirements of these scenarios. For theoretical accuracy benchmark comparison scenarios, the Baldeaux exact simulation method may be used, but strictly under benchmark parameter environments. The remaining terminal simulation methods are not recommended for use in practical pricing scenarios and should only be utilized for specific academic comparative analyses.

\section{Research Limitations}

This thesis has conducted systematic theoretical and empirical research centered on the option pricing problem under the 3/2 stochastic volatility model, forming a relatively complete analytical framework. However, constrained by the research scope, data conditions, and model specifications, certain limitations remain, primarily reflected in the following four aspects:

First, there is room for expansion in the dimensions of model specification. This thesis focuses on the 3/2 stochastic volatility model under a pure diffusion framework and has not introduced market friction factors such as underlying price jump terms, stochastic interest rates, or stochastic liquidity into the model. Although the pure diffusion model can already fit the volatility smile morphology throughout the full cycle of crash events quite well, the jump characteristics of underlying prices and instantaneous changes in market liquidity during extreme conditions may still exert a systematic impact on option pricing. There is still further room to enrich the model specification to portray real market structures.

Second, the coverage scope of empirical research needs to be broadened. The empirical analysis in this thesis is based solely on BTC option data from the Deribit exchange, without covering other mainstream crypto assets like ETH, nor extending to underlying assets in traditional financial markets such as equities, commodities, or foreign exchange. Significant differences exist in the volatility characteristics, market microstructures, and investor structures across different underlying assets. The adaptability and parameter characteristics of the 3/2 model across various markets still require cross-validation using multi-dimensional empirical data.

Third, there are limitations in the evaluation scope of pricing methods. The performance evaluation of pricing methods in this thesis focuses solely on European options and does not conduct analyses on mainstream path-dependent derivatives such as American options, barrier options, and Asian options. Method construction, accuracy optimization, and performance evaluation for pricing path-dependent options under the 3/2 model remain inadequately covered by existing research and represent a core limitation of the research scope of this thesis.

Fourth, there is room for improvement in the richness of the parameter calibration framework. The Medvedev-Scaillet asymptotic expansion method adopted in this thesis is more suitable for short-maturity options, and there is still room to optimize calibration accuracy for ultra-long-maturity options. Concurrently, this thesis only employs a single calibration method for parameter estimation and has not introduced methods like characteristic function-based global calibration or Bayesian calibration to conduct robust comparisons of parameter estimation results. The diversity and robustness of the calibration framework can still be further enhanced.

\section{Future Research Outlook}

Based on the research conclusions and existing limitations of this thesis, further in-depth research can be conducted in the following five directions in the future to perfect the pricing theory and practical application system of the 3/2 stochastic volatility model:

First, expand the theoretical specifications of the 3/2 stochastic volatility model. Building upon the pure diffusion model, factors such as synchronous jumps, independent jumps, stochastic interest rates, and stochastic liquidity can be introduced to construct an extended 3/2 model that better aligns with real market structures. Derive the closed-form solutions of characteristic functions and option pricing formulas for the corresponding extended models, quantify the marginal impact of different expansion factors on option pricing, and further enhance the model's adaptability to extreme market environments. Concurrently, conduct horizontal comparisons of pricing accuracy between the extended models and the benchmark model to clearly define the applicable scenarios of different model setups.

Second, expand the applicable scope of pricing methods and optimize computational performance. Targeting mainstream path-dependent derivatives such as American options, barrier options, and Asian options, construct highly efficient pricing methods adapted to the variance process characteristics of the 3/2 model. Complete systematic evaluations of the pricing accuracy, computational efficiency, and numerical stability of different methods to form a 3/2 model pricing system covering a full range of derivative products. Meanwhile, by integrating technologies such as GPU parallel computing and vectorization optimization, further elevate the computational efficiency of large-scale Monte Carlo simulations and frequency-domain methods to meet the high-frequency real-time pricing demands in actual business operations.

Third, expand the market coverage dimensions of empirical research. Extend the empirical analysis from BTC options to other mainstream crypto assets like ETH, as well as to traditional financial sectors such as global equity markets, commodity markets, and foreign exchange markets. Validate the adaptability of the 3/2 model across different volatility characteristics and market structures to clearly delineate the model's universal application boundaries. Concurrently, based on longer-cycle historical data, conduct large-sample horizontal comparisons between the 3/2 model and mainstream stochastic volatility models like the Heston model and SABR model. Quantify the differences in pricing accuracy and market explanatory power among various models to provide empirical evidence for model selection across different markets.

Fourth, optimize the parameter calibration framework of the 3/2 model. Addressing the shortcomings of the asymptotic expansion method in calibrating long-term options, introduce methods such as characteristic function-based global calibration, Bayesian calibration, and machine learning-assisted calibration to elevate parameter calibration accuracy and stability in extreme market scenarios. Simultaneously, construct a real-time dynamic calibration system for model parameters to achieve high-frequency updates of model parameters under extreme conditions, adapting to the real-time risk management demands of derivatives markets.

Fifth, construct a full-process risk management system for derivatives based on the 3/2 model. Building on pricing research, derive closed-form solutions and efficient numerical computation methods for option Greeks under the 3/2 model. Establish model-based tail risk measurement systems such as VaR and CVaR to quantify the risk exposure of option portfolios during extreme conditions. Concurrently, based on the time-varying characteristics of model parameters, build measurement models for volatility risk premiums, providing theoretical support and empirical evidence for hedging strategy design and dynamic asset allocation in the crypto derivatives market.